\documentclass[12pt]{article}
\usepackage[utf8]{inputenc}
\usepackage[T1]{fontenc}
\usepackage[margin=1in]{geometry}
\usepackage{setspace}
\usepackage{xcolor}
\usepackage[colorlinks=true, urlcolor=blue!60!black, linkcolor=blue!60!black]{hyperref}
\setlength{\parindent}{2em}
\setlength{\parskip}{1.25em}
\renewcommand{\baselinestretch}{1.0}
\begin{document}
\begin{center}
{\LARGE \textit{Griswold v. Connecticut}, 381 U.S. 479 (1965)\par}
\vspace{0.8em}
{\large Case Brief --- Four-Page Edition\par}
\end{center}
\vspace{1em}

\section{Memory Jogger}

The Court invalidated Connecticut's criminal ban on contraceptive use by married couples, recognizing a constitutional right to marital privacy derived from ``penumbras, formed by emanations'' of specific guarantees in the Bill of Rights (\href{https://supreme.justia.com/cases/federal/us/381/479/}{381 U.S. 479}).

\section{Detailed Case Facts}

Since 1879 Connecticut had made it a crime for ``[a]ny person'' to use any drug or instrument ``for the purpose of preventing conception,'' and its general accessory statute allowed anyone who assisted or counseled an offense to be punished as a principal. Earlier challenges had failed on justiciability grounds: \textit{Tileston v. Ullman} (1943) for lack of standing and \textit{Poe v. Ullman} (1961) as unripe because the law was rarely enforced.

To force a test case, Estelle Griswold, executive director of the Planned Parenthood League of Connecticut, and Dr. C. Lee Buxton, its medical director and a Yale professor, opened a birth control clinic in New Haven. From November 1 to 10, 1961, they gave married persons information and medical advice on preventing conception, examined wives, and prescribed contraceptives. They were arrested and charged as accessories.

\section{Procedural History}

After a bench trial, Griswold and Buxton were convicted as accessories and fined \$100 each. The Appellate Division and the Connecticut Supreme Court of Errors affirmed. The U.S. Supreme Court noted probable jurisdiction and reversed on June 7, 1965.

\section{Judicial Votes}

Reversed, \textbf{7--2}.
\begin{itemize}
\item \textbf{Opinion of the Court:} Justice Douglas, joined by Chief Justice Warren and Justices Clark, Brennan, and Goldberg.
\item \textbf{Concurrence:} Justice Goldberg, joined by Chief Justice Warren and Justice Brennan.
\item \textbf{Concurring in the judgment:} Justice Harlan; Justice White.
\item \textbf{Dissents:} Justice Black, joined by Justice Stewart; Justice Stewart, joined by Justice Black.
\end{itemize}

\section{Holding}

Accessories convicted for helping married couples have standing to assert the couples' rights. A state law forbidding married couples to use contraceptives violates a fundamental right of marital privacy protected within the ``penumbras'' of specific Bill of Rights guarantees, applied to the States through the Fourteenth Amendment; by forbidding \textit{use}, the law sweeps unnecessarily broadly. Douglas's rhetoric about marriage, his suggestion that States might regulate ``manufacture or sale,'' and Goldberg's assurance that the decision does not disturb laws against adultery and fornication were dicta.

\section{Analysis of Opinions}

\textbf{Douglas (Court).} The defendants' convictions distinguished \textit{Tileston}. Declining to revive \textit{Lochner} or sit as a ``super-legislature,'' Douglas found that the First, Third, Fourth, Fifth, and Ninth Amendments create ``zones of privacy.'' A ban on use had a ``maximum destructive impact'' on marriage, raising the repulsive prospect of police searching marital bedrooms.

\textbf{Goldberg (with Warren and Brennan).} The Ninth Amendment, though not an independent source of rights, shows that the Framers believed fundamental liberties exist beyond those enumerated; judges should consult the ``traditions and [collective] conscience of our people.'' Connecticut could punish adultery directly, so its law failed strict scrutiny.

\textbf{Harlan and White (in the judgment).} Harlan, relying on his \textit{Poe} dissent, grounded the right directly in the Due Process Clause, which protects values ``implicit in the concept of ordered liberty,'' with restraint coming from history and tradition. White found no rational connection between banning married couples' use and deterring illicit affairs.

\textbf{Dissents (Black and Stewart).} Black found the law offensive but unconstitutional only if a specific provision forbade it, and attacked the concurrences as \textit{Lochner}-style ``natural law due process'' making the Court a ``day-to-day constitutional convention.'' Stewart called it ``an uncommonly silly law,'' walked through each cited amendment to show none was violated, and left repeal to the legislature.

\section{Examples: Future Applications}

\begin{enumerate}
\item \textbf{Same-Side (Invalid): Ban on Distribution to Unmarried Adults.} Restricting the privacy interest in whether to bear or beget a child by marital status fails, as \textit{Eisenstadt v. Baird} later held.
\item \textbf{Same-Side (Invalid): Two-Child Sterilization Mandate.} Requiring married couples to be sterilized after two children invades marital intimacy, as Goldberg noted, and fails strict scrutiny.
\item \textbf{Opposite-Side (Valid): Prescription Requirement.} Requiring a prescription and pharmacy dispensing for hormonal contraceptives regulates distribution and safety, not use, tracking Douglas's distinction.
\item \textbf{Opposite-Side (Valid): Domestic Abuse in the Home.} A spouse prosecuted for marital rape cannot invoke \textit{Griswold}, which protects consensual intimacy, not coercion.
\item \textbf{Fence-Sitter: Parental Notification for Minors.} A rule requiring parental notice before prescribing birth control to minors under 16 pits the minor's privacy against the State's interest and parents' rights under \textit{Meyer} and \textit{Pierce}.
\end{enumerate}

\section{Critique}

\textbf{Originalist and textualist critiques.} Robert Bork argued that Douglas leapt without justification from specific guarantees to a freestanding ``right of privacy,'' leaving no principled limit, and Justice Thomas in \textit{Dobbs} urged reconsidering substantive due process precedents, naming \textit{Griswold}.

\textbf{Progressive and feminist critiques.} John Hart Ely doubted unenumerated substantive due process rights. Feminist scholars note the ``paradoxes of privacy'': a marital frame excluded unmarried women and risked shielding domestic subordination. Later Courts largely abandoned penumbras for Harlan's due process approach.

\section{Key Quotations}

\begin{quote}
``The foregoing cases suggest that specific guarantees in the Bill of Rights have penumbras, formed by emanations from those guarantees that help give them life and substance.''
\end{quote}
\noindent --- Justice Douglas, Opinion of the Court, 381 U.S. at 484.

\begin{quote}
``Would we allow the police to search the sacred precincts of marital bedrooms for telltale signs of the use of contraceptives? The very idea is repulsive to the notions of privacy surrounding the marriage relationship.''
\end{quote}
\noindent --- Justice Douglas, Opinion of the Court, 381 U.S. at 485--486.

\begin{quote}
``I think this is an uncommonly silly law. ... But we are not asked in this case to say whether we think this law is unwise, or even asinine. We are asked to hold that it violates the United States Constitution. And that I cannot do.''
\end{quote}
\noindent --- Justice Stewart, dissenting, 381 U.S. at 527.

\end{document}
